\chapter{Protected Sectors and Superconformal Indices}
\label{ch:volume-iii-protected-sectors-superconformal-indices}

\section{Cohomological Protected Data}

Let the theory be a unitary superconformal field theory on
\(\mathbb R^D\).  Radial quantization gives a Hilbert space
\(\Hilb(S^{D-1})\), and the superconformal algebra acts by densely defined
operators on a common invariant domain.  Choose a fermionic supercharge
\[
  Q
\]
and let \(Q^\dagger\) be its adjoint in radial quantization.  Suppose
\[
  Q^2=0,
  \qquad
  \mathcal B=\{Q,Q^\dagger\}
  \label{eq:bps-positive-operator}
\]
is a nonnegative self-adjoint operator.  In examples, \(\mathcal B\) is a
linear combination of the dilatation generator, rotation generators, and
R-symmetry generators.  A state \(\ket\psi\) is \(Q\)-BPS when
\[
  \mathcal B\ket\psi=0.
\]
Since
\[
  \braket{\psi}{\mathcal B\psi}
  =
  \|Q\ket\psi\|^2+\|Q^\dagger\ket\psi\|^2,
\]
the BPS condition is equivalent to
\[
  Q\ket\psi=0,
  \qquad
  Q^\dagger\ket\psi=0.
\]

The \(Q\)-cohomology of states is
\begin{equation}
  H_Q(\Hilb)
  =
  \frac{\ker Q}{\operatorname{im} Q}.
  \label{eq:q-cohomology-hilbert-space}
\end{equation}
Under standard domain assumptions for the Hodge decomposition associated with
\(Q\) and \(Q^\dagger\), every cohomology class has a unique harmonic
representative annihilated by both \(Q\) and \(Q^\dagger\).  Thus
\[
  H_Q(\Hilb)\simeq \ker\mathcal B .
\]
The protected sector is the part of the Hilbert space, or equivalently the
part of the local operator spectrum by the state-operator correspondence, that
survives in this cohomology.

\begin{figure}[H]
  \centering
  \resizebox{0.90\textwidth}{!}{%
  \begin{tikzpicture}[x=1cm,y=1cm,every node/.style={font=\scriptsize}]
    \tikzset{
      arrow/.style={-{Latex[length=2mm]},line width=0.85pt},
      box/.style={draw,rounded corners=4pt,line width=0.8pt,fill=blue!4,
        minimum width=3.0cm,minimum height=0.86cm,align=center},
      harmonic/.style={draw,circle,line width=0.9pt,fill=green!8,
        minimum size=1.02cm,align=center}
    }
    \node[box] (hilb) at (-4.3,0.95) {\(\Hilb(S^{D-1})\)};
    \node[box] (ker) at (0,0.95) {\(\ker Q\)};
    \node[box] (im) at (0,-0.95) {\(\operatorname{im}Q\)};
    \node[harmonic] (harm) at (4.25,0.95) {\(\ker\mathcal B\)};
    \draw[arrow] (hilb)--(ker) node[midway,above] {\(Q\ket\psi=0\)};
    \draw[arrow] (im)--(ker) node[midway,right] {quotient};
    \draw[arrow] (ker)--(harm);
    \node[align=center] at (2.85,1.52) {harmonic\\representative};
    \node[green!45!black,align=center] at (0,-2.0)
      {\(\mathcal B=\{Q,Q^\dagger\}\) selects the protected cohomology};
  \end{tikzpicture}
  }
  \caption{A protected sector is the \(Q\)-cohomology represented by states
  annihilated by the positive operator \(\mathcal B=\{Q,Q^\dagger\}\).}
  \label{fig:protected-sector-cohomology}
\end{figure}

\section{Short Multiplets and Recombination}

Let \(\mathfrak g_{\rm sc}\) be the superconformal algebra and let
\(\mathfrak h\) be a Cartan subalgebra containing the dilatation generator
\(D_{\rm cyl}\) on the cylinder, spatial rotations, and R-symmetry
generators.  A superconformal primary is a state annihilated by the conformal
raising operators and by the superconformal raising operators.  Acting with
Poincare supercharges generates a supermultiplet.

Unitarity gives inequalities of the form
\begin{equation}
  \Delta\ge f(\rho,j,R),
  \label{eq:superconformal-shortening-bound}
\end{equation}
where \(\Delta\) is the scaling dimension, \(\rho\) denotes Lorentz spin
data, \(j\) denotes possible little-group labels, and \(R\) denotes
R-symmetry weights.  A short multiplet is a multiplet for which at least one
unitarity inequality is saturated.  Saturation creates null descendants and
therefore a reduced representation.

Along a conformal manifold, short multiplets can combine into long multiplets
when quantum numbers allow a recombination rule of the schematic form
\[
  L_{\Delta\to\Delta_*}
  \longrightarrow
  S_1\oplus S_2\oplus\cdots .
\]
Here \(L\) is a long multiplet whose dimension approaches a threshold
\(\Delta_*\), and the \(S_i\) are short multiplets at threshold.  An individual
short multiplet can therefore be protected by charges and selection rules, or
it can participate in a recombination pattern.  A protected index records the
net cohomological information stable under such pairings.

\section{The Superconformal Index}

Let \(F\) be fermion number, and let \(Q_a\) be mutually commuting bosonic
charges that commute with \(Q\), \(Q^\dagger\), and \(\mathcal B\).  Introduce
fugacities \(y_a\).  The superconformal index is the graded trace
\begin{equation}
  \mathcal I(y_a)
  =
  \operatorname{Tr}_{\Hilb(S^{D-1})}
  \left[
    (-1)^F
    \ee^{-\beta\mathcal B}
    \prod_a y_a^{Q_a}
  \right],
  \label{eq:superconformal-index-definition}
\end{equation}
with \(\beta>0\).  States with \(\mathcal B>0\) occur in \(Q\)-pairs
\[
  \ket\psi,\qquad Q\ket\psi,
\]
with opposite fermion number and the same commuting charges.  Their
contributions to \eqref{eq:superconformal-index-definition} cancel.  Hence
\[
  \frac{\partial\mathcal I}{\partial\beta}=0,
  \qquad
  \mathcal I(y_a)
  =
  \operatorname{Tr}_{H_Q(\Hilb)}
  \left[
    (-1)^F\prod_a y_a^{Q_a}
  \right].
\]
The index is therefore a protected function of fugacities.  It is insensitive
to continuous deformations that preserve the supercharge, the commuting
charges, the trace-class property of the graded trace, and the absence of
states escaping through a continuum at \(\mathcal B=0\).

\begin{figure}[H]
  \centering
  \resizebox{0.88\textwidth}{!}{%
  \begin{tikzpicture}[x=1cm,y=1cm,every node/.style={font=\scriptsize}]
    \tikzset{
      arrow/.style={-{Latex[length=2mm]},line width=0.85pt},
      state/.style={draw,rounded corners=3pt,fill=blue!4,inner sep=4pt,
        align=center},
      bps/.style={draw,rounded corners=3pt,fill=green!8,inner sep=4pt,
        align=center}
    }
    \node[state] (a) at (-3.7,0.70) {\(\ket\psi\)\\\(\mathcal B>0\)};
    \node[state] (b) at (-1.0,0.70) {\(Q\ket\psi\)\\\(\mathcal B>0\)};
    \draw[arrow] (a)--(b) node[midway,above] {\(Q\)};
    \node[orange!85!black] at (-2.35,-0.10) {opposite \((-1)^F\)};
    \draw[orange!85!black,line width=0.85pt] (-4.35,-0.35)--(-0.35,-0.35);
    \node[orange!85!black] at (-2.35,-0.72) {paired contribution cancels};
    \node[bps] (c) at (3.0,0.45) {\(\mathcal B=0\)\\cohomology class};
    \draw[arrow,green!45!black] (0.15,0.30)--(1.85,0.42);
    \node[green!45!black,align=center] at (3.0,-0.82)
      {index receives only\\protected contributions};
  \end{tikzpicture}
  }
  \caption{The graded trace pairs non-BPS states by the action of \(Q\), while
  cohomology classes at \(\mathcal B=0\) remain.}
  \label{fig:index-pairing-cancellation}
\end{figure}

\section{Protected Operator Products}

Let \(\mathcal O_1\) and \(\mathcal O_2\) be local operators satisfying
\[
  [Q,\mathcal O_i\}=0,
\]
where the bracket is graded.  Their product inside a correlation function is
also \(Q\)-closed:
\[
  [Q,\mathcal O_1(x)\mathcal O_2(0)\}=0.
\]
If \(\mathcal O_i\) is shifted by a \(Q\)-exact operator, the product changes
by a \(Q\)-exact operator.  Therefore the OPE descends to a product on
cohomology:
\begin{equation}
  [\mathcal O_i]\,[\mathcal O_j]
  =
  \sum_k C_{ij}{}^k\,[\mathcal O_k],
  \label{eq:protected-ope-cohomology-product}
\end{equation}
where brackets denote \(Q\)-cohomology classes.  This protected product can
be finite-dimensional, as in a chiral ring, or infinite-dimensional, as in a
protected chiral algebra.

For a four-dimensional \(\mathcal N=2\) superconformal theory, one can choose
a nilpotent supercharge \(Q_{\rm Schur}\) and restrict local operators to a
fixed two-plane with complex coordinate \(z\).  In the corresponding
cohomology, translations in \(\bar z\) and in the transverse directions are
\(Q_{\rm Schur}\)-exact.  The OPE of cohomology classes is holomorphic:
\[
  [\mathcal O_i](z)[\mathcal O_j](0)
  =
  \sum_k f_{ij}{}^k(z)[\mathcal O_k](0).
\]
The resulting two-dimensional chiral algebra is protected data of the
four-dimensional theory.  With conventional normalizations,
\[
  c_{\rm 2d}=-12\,c_{\rm 4d},
  \qquad
  k_{\rm 2d}=-\frac12\,k_{\rm 4d},
\]
relating the chiral-algebra central charge and affine levels to the
four-dimensional stress-tensor and flavor-current two-point coefficients.

\begin{figure}[H]
  \centering
  \resizebox{0.92\textwidth}{!}{%
  \begin{tikzpicture}[x=1cm,y=1cm,every node/.style={font=\scriptsize}]
    \tikzset{
      arrow/.style={-{Latex[length=2mm]},line width=0.85pt},
      plane/.style={fill=blue!5,draw=blue!45!black,line width=0.8pt},
      op/.style={circle,fill=violet!80!black,inner sep=2pt}
    }
    \draw[plane] (-4.3,-0.95) rectangle (0.4,1.05);
    \node[blue!65!black] at (-1.95,1.38) {chosen two-plane};
    \draw[->,thick] (-3.95,-0.65)--(-2.95,-0.65) node[right] {\(z\)};
    \draw[->,thick] (-3.95,-0.65)--(-3.95,0.35) node[above] {\(\bar z\)};
    \node[op] at (-2.85,0.15) {};
    \node[op] at (-1.55,0.34) {};
    \node[violet!80!black] at (-2.82,0.55) {\([\mathcal O_i]\)};
    \node[violet!80!black] at (-1.52,0.74) {\([\mathcal O_j]\)};
    \draw[arrow] (0.85,0.05)--(2.10,0.05);
    \node[align=center] at (1.48,0.55) {\(Q\)-cohomology\\OPE};
    \node[draw,rounded corners=4pt,fill=green!8,align=center,
      minimum width=2.9cm,minimum height=1.0cm] at (3.75,0.05)
      {protected\\chiral algebra};
    \node[orange!85!black,align=center] at (-1.95,-1.45)
      {\(\bar z\)-motion and transverse motion are \(Q\)-exact};
  \end{tikzpicture}
  }
  \caption{A protected chiral algebra arises when a \(Q\)-cohomology of local
  operators has a holomorphic OPE on a distinguished two-plane.}
  \label{fig:protected-chiral-algebra-plane}
\end{figure}

\section{Protected Data and Full Data}

Protected sectors are part of the complete conformal data.  They supply exact
dimensions, exact OPE coefficients, graded traces, and sometimes exact
subalgebras.  They constrain the remaining dynamical data through ordinary
crossing, through superconformal Ward identities, and through duality on
conformal manifolds.

The relation is one-directional as a matter of information: the protected
sector gives exact projections of the full theory, while the full theory
contains long multiplets and unprotected OPE coefficients in addition to the
protected projection.  A systematic analysis therefore uses protected data as
anchors inside the larger crossing problem rather than as a replacement for
the full operator algebra.
